\documentclass[11pt,a4paper]{article}
\usepackage[T1]{fontenc}
\usepackage[utf8]{inputenc}
\usepackage{lmodern}
\usepackage[margin=27mm,headheight=14pt]{geometry}
\usepackage{amsmath,amssymb,amsthm,mathtools}
\usepackage{microtype}
\usepackage{booktabs,array,longtable}
\usepackage{enumitem}
\usepackage{fancyhdr}
\usepackage{hyperref}
\usepackage{bookmark}
\hypersetup{hidelinks,pdftitle={A non-Baire translation-invariant ideal with a perfect translation-almost-disjoint family},pdfauthor={ChatGPT},pdfsubject={A proposed affirmative solution to Question 1 of Filipow and Tryba}}
\pagestyle{fancy}
\fancyhf{}
\fancyhead[L]{\small Non-Baire ideals and translated almost disjointness}
\fancyhead[R]{\small Research note}
\fancyfoot[C]{\thepage}
\renewcommand{\headrulewidth}{0.3pt}
\setlength{\parindent}{1.2em}
\setlength{\parskip}{3pt}
\setlength{\emergencystretch}{2em}
\allowdisplaybreaks[2]
\newtheorem{theorem}{Theorem}[section]
\newtheorem{lemma}[theorem]{Lemma}
\newtheorem{proposition}[theorem]{Proposition}
\newtheorem{corollary}[theorem]{Corollary}
\theoremstyle{definition}
\newtheorem{definition}[theorem]{Definition}
\newtheorem{example}[theorem]{Example}
\newtheorem{problem}[theorem]{Problem}
\theoremstyle{remark}
\newtheorem{remark}[theorem]{Remark}
\newcommand{\N}{\omega}
\newcommand{\Z}{\mathbb Z}
\newcommand{\R}{\mathbb R}
\newcommand{\I}{\mathcal I}
\newcommand{\J}{\mathcal J}
\newcommand{\U}{\mathcal U}
\newcommand{\M}{\mathcal M}
\newcommand{\A}{\mathcal A}
\newcommand{\Pow}{\mathcal P}
\newcommand{\Fin}{\mathrm{Fin}}
\newcommand{\cfr}{\mathfrak c}
\newcommand{\supp}{\sigma}
\newcommand{\sh}[2]{\tau_{#1}(#2)}
\newcommand{\eqfin}{=^{*}}
\newcommand{\subfin}{\subseteq^{*}}
\newcommand{\bp}{\mathrm{BP}}
\newcommand{\restr}{\mathbin{\upharpoonright}}
\DeclareMathOperator{\val}{val}
\title{\vspace{-1.5em}\textbf{A non-Baire translation-invariant ideal\protect\\
with a perfect translation-almost-disjoint family}\protect\\[0.65em]
\large An affirmative construction for a question of Filip\'ow and Tryba}
\author{ChatGPT\thanks{Prepared in response to a request for a randomly selected open-problem investigation. This is an unrefereed research note, not an independently certified or formally verified result.}}
\date{September 20, 2026}

\begin{document}
\maketitle
\begin{abstract}
We give a ZFC construction intended to answer Question~1 of Filip\'ow and
Tryba, \emph{Densities for sets of natural numbers vanishing on a given
family}, affirmatively. For every free ultrafilter $\U$ on $\omega$, an
explicit partition into finite intervals defines a proper
translation-invariant ideal $\I_\U$ without the Baire property. There is a
continuous injection $x\mapsto A_x$ from $2^\omega$ into $\I_\U^+$ such
that $A_x\cap(A_y+k)$ is finite whenever $x\ne y$ and $k\in\mathbb Z$.
Thus the positive family can be perfect, and its translated intersections
can be genuinely finite rather than merely ideal-small. The construction
combines ultrafilter block supports with widely separated labels on a
binary tree. A direct fusion argument proves the category assertion
without relying on a preservation claim about arbitrary continuous
preimages. We also establish quantitative intersection bounds, maximal
ultrafilter traces, a rich maxitive abstract upper density with this exact
null ideal, and an extension to countable permutation actions. Exact
finite tests accompany the proof; they check its explicit arithmetic and
fusion invariants, not the infinitary assertions. The literature search
found the question explicitly posed in the original source, but does not
constitute a definitive priority determination.
\end{abstract}

\vspace{0.3em}
\noindent\textbf{Keywords.} Ideals on the integers; Baire property; free
ultrafilters; translation almost disjointness; abstract upper densities.

\newpage
\tableofcontents
\newpage

\section{The problem, the result, and its status}
\label{sec:intro}

\subsection{Random selection and problem selection}
The mathematical area was chosen before the problem. A single call to a
pseudorandom generator, seeded with 128 bits from the operating system,
selected index~22 from a fixed, zero-indexed list of 24 mathematical
areas. The result was \emph{set theory}. The seed, the complete list, and
reproduction instructions appear in Appendix~\ref{app:random} and in
\texttt{random\_selection.json}. There was no reroll of the area.

Within set theory, ideals on a countable set offered a useful balance:
the motivating question is genuinely infinitary, but potential witnesses
can be built from finite blocks, binary trees, and ultrafilters. The
selected problem is Question~1 of Filip\'ow and Tryba~\cite{FT}. In the
notation developed below, its content is the following.

\begin{problem}[Filip\'ow--Tryba, Question 1]
\label{prob:FT}
Find a translation-invariant proper ideal $\I$ on $\omega$ for which
$\I$ lacks the Baire property and $\I^+$ contains a
translation-almost-disjoint family of size $\cfr=2^{\aleph_0}$.
\end{problem}

The original uses the positive integers. Passing between $\omega$ and
the positive integers by $n\mapsto n+1$ does not change this question.
The translation conventions and the harmless finite boundary effects
are made explicit in Section~\ref{sec:prelim}.

\subsection{Main statement}
\begin{theorem}[Main construction]
\label{thm:main}
In ZFC, for every free ultrafilter $\U$ on $\omega$ there exist a proper
ideal $\I_\U\subseteq\Pow(\omega)$ and a continuous injection
\[
             \Phi:2^\omega\longrightarrow 2^\omega,
             \qquad \Phi(x)=A_x,
\]
with the following properties:
\begin{enumerate}[label=\textup{(\roman*)},itemsep=2pt]
\item $\I_\U$ contains $\Fin$, is translation invariant, and does not
have the Baire property.
\item Every $A_x$ is $\I_\U$-positive.
\item If $x\ne y$ and $k\in\mathbb Z$, then
$A_x\cap(A_y+k)$ is finite. If $k\ne0$, this conclusion also holds for
$x=y$.
\item $\I_\U$ is tall. More strongly, its restriction to any positive
set does not have the Baire property in the corresponding Cantor space.
\end{enumerate}
The image of $\Phi$ is a compact perfect family of cardinality $\cfr$.
Both the block partition and the map $\Phi$ are independent of $\U$.
\end{theorem}

The proof is given in Sections~\ref{sec:category}--\ref{sec:tree}.
No additional hypothesis such as the continuum hypothesis is used. The
construction is explicit \emph{relative to a free ultrafilter}; this is
not a claim that membership in the ideal is algorithmically decidable.

The theorem answers the existential question by strengthening its
witness. An ideal-small intersection is allowed in the problem, whereas
our intersections are finite. A family of continuum cardinality is
requested, whereas our family is the continuous image of Cantor space.
These stronger properties will also make it possible to prescribe
positive values of an associated abstract density.

\subsection{What is, and is not, being claimed}
The original article was published in 2020; its arXiv version was posted
in 2023. Question~1 is visible on page~7 of the ten-page preprint~\cite{FT}.
A targeted search on September~20, 2026 checked this source, the author's
publication and citation lists, and related accessible later papers,
including~\cite{KL,Path}. No resolution of this particular question was
located. These checks do not exclude an unindexed paper, an unpublished
argument, or another treatment under different terminology. A search
record is included in the archive.

The mathematical contribution claimed here is the complete written
construction and proof, not a verified assertion of priority. The note
has not been refereed or checked in a proof assistant. The explicit
formulas and several finite versions of the combinatorial invariants
were checked by an executable program. This provides protection against
indexing errors but is not a substitute for the proofs.

The category criterion used below is classical, and the maximal-family
method for constructing abstract densities is established in the
literature~\cite{Talagrand,DNJ,FT}. These ingredients are not presented as
new. The later questions in~\cite{FT} about all invariant ideals or
translation-maximal ideals are not answered by this construction.

\section{Preliminaries and translation conventions}
\label{sec:prelim}

Write $\omega=\{0,1,2,\ldots\}$ and identify a subset of a countable set
$X$ with its characteristic function in $2^X$. This endows $\Pow(X)$ with
the product topology. After any enumeration of $X$, this is Cantor space.
A subset of this space has the \emph{Baire property} if its symmetric
difference with an open set is meager. A meager set is a countable union
of nowhere dense sets.

\begin{definition}
An \emph{ideal} on $X$ is a family $\I\subseteq\Pow(X)$ that contains all
finite subsets of $X$, is hereditary, is closed under finite unions, and
does not contain $X$. Its positive sets are
$\I^+=\Pow(X)\setminus\I$. It is \emph{tall} if every infinite
$D\subseteq X$ has an infinite subset belonging to $\I$.
\end{definition}

Every ideal is invariant under finite symmetric differences. Indeed,
if $A\in\I$ and $F$ is finite, then $A\mathbin\triangle F\subseteq
A\cup F\in\I$; the converse follows by applying the same observation
again. We write $A\eqfin B$ when $A\mathbin\triangle B$ is finite and
$A\subfin B$ when $A\setminus B$ is finite.

For $k\in\mathbb Z$ define the truncated shift
\begin{equation}
\label{eq:shift}
                 \sh{k}{A}=(A+k)\cap\omega.
\end{equation}
It preserves unions and inclusions. For all integers $h,k$,
\begin{equation}
\label{eq:composition}
           \sh{k}{\sh{h}{A}}\eqfin\sh{k+h}{A}.
\end{equation}
To check this, the only possible loss is caused by an element shifted
below zero at the intermediate step and shifted back into $\omega$ at
the next step. Such original elements lie in a fixed finite initial
segment, depending on $h$ and $k$. An ideal is \emph{translation
invariant} if $A\in\I$ is equivalent to $\sh{k}{A}\in\I$ for every $k$.

\begin{definition}
A family $\A\subseteq\I^+$ is \emph{$\I$-translation-almost-disjoint},
or \emph{$\I$-TAD}, if
\[
             A\cap\sh{k}{B}\in\I
     \quad(A,B\in\A,\ A\ne B,\ k\in\mathbb Z).
\]
We call the stronger property that all these intersections are finite
\emph{finite translation almost disjointness}.
\end{definition}

A free ultrafilter $\U$ on $\omega$ is a maximal proper filter containing
all cofinite sets and no finite sets. In particular, for every
$S\subseteq\omega$ exactly one of $S$ and $\omega\setminus S$ belongs to
$\U$. Its dual ideal is
\[
           \U^{\mathrm d}=\{S\subseteq\omega:S\notin\U\}.
\]
For example, if $S,T\notin\U$, their complements belong to $\U$, so
$\omega\setminus(S\cup T)\in\U$ and $S\cup T\notin\U$.
Membership in $\U$ is unaffected by finite changes.

Free ultrafilters exist in ZFC by extending the cofinite filter to an
ultrafilter. No special ultrafilter, such as a selective ultrafilter or a
$P$-point, will be needed. In all that follows, $\U$ is an arbitrary
fixed free ultrafilter.

\section{An elementary category criterion for ideals}
\label{sec:category}

The next criterion is the form of the classical category
characterization of ideals needed here; compare
Talagrand~\cite{Talagrand} and its formulation in~\cite[Theorem~1.7]{FT}.
We include the proof to make the category step inspectable.

\begin{lemma}
\label{lem:zeroone}
If a proper ideal on a countably infinite set has the Baire property,
then it is meager.
\end{lemma}
\begin{proof}
Enumerate the ground set as $\omega$. Suppose that $\I$ has the Baire
property and is not meager. There is a nonempty basic cylinder $[s]$ in
which $\I$ is comeager. For every binary word $t$ of length $|s|$, flipping
a suitable subset of the first $|s|$ coordinates gives a homeomorphism
from $[s]$ onto $[t]$. It preserves $\I$ because $\I$ is invariant under
finite changes. Thus $\I$ is comeager in every such cylinder, and hence
in all of $2^\omega$.

The complementation map $A\mapsto\omega\setminus A$ is a
homeomorphism. Its image of $\I$ is therefore also comeager. This image
is disjoint from $\I$: otherwise an ideal would contain both a set and
its complement, and so would contain $\omega$. Two comeager subsets of
Cantor space cannot be disjoint, by the Baire category theorem. This is
a contradiction.
\end{proof}

\begin{lemma}[Finite-block category criterion]
\label{lem:criterion}
A proper ideal $\I$ on $\omega$ has the Baire property if and only if
there are integers
\[
               0=t_0<t_1<t_2<\cdots
\]
such that no member of $\I$ contains infinitely many whole intervals
$[t_n,t_{n+1})$.
\end{lemma}
\begin{proof}
Suppose first that $\I$ has the Baire property. By
Lemma~\ref{lem:zeroone}, it is contained in $\bigcup_j K_j$, where each
$K_j\subseteq2^\omega$ is closed and nowhere dense.

We recursively choose $t_{n+1}>t_n$ and a binary word $u_n$ on the
coordinates $[t_n,t_{n+1})$ so that
\begin{equation}
\label{eq:pattern}
 [v\,u_n]\cap(K_0\cup\cdots\cup K_n)=\varnothing
             \qquad\text{for every }v\in2^{t_n}.
\end{equation}
Here $v\,u_n$ denotes concatenation. There are only finitely many such
prefixes $v$. Start with an empty proposed common suffix. For the first
prefix, extend the suffix until the resulting cylinder avoids the
finite union of the $K_j$'s. For the next prefix, extend the same suffix
further to obtain avoidance for that prefix as well. Previously obtained
avoidance persists under extension. After processing every prefix, add
a coordinate if necessary to make the suffix nonempty. Now set
$t_{n+1}$ equal to the resulting total length. Nowhere density justifies
every extension.

If a point $z\in2^\omega$ agrees with $u_n$ on infinitely many of these
intervals, it belongs to none of the $K_j$. For any fixed $j$, choose a
matching $n\ge j$ and apply~\eqref{eq:pattern} to the actual prefix of
$z$.

Suppose, toward a contradiction, that some $A\in\I$ contains infinitely
many full intervals $[t_n,t_{n+1})$. On those intervals choose a subset
$D\subseteq A$ whose characteristic function is $u_n$, and put no
points into $D$ elsewhere. Then $D$ avoids every $K_j$, so $D\notin\I$.
This contradicts heredity.

Conversely, suppose the stated partition exists. For each $m$, the set
\[
 E_m=\{A\subseteq\omega:
           [t_n,t_{n+1})\nsubseteq A\text{ for every }n\ge m\}
\]
is closed and nowhere dense. It is closed because each single-block
condition is clopen. It has empty interior because any basic cylinder
can be extended by putting an entire sufficiently late block into the
set. The hypothesis gives $\I\subseteq\bigcup_m E_m$. Thus $\I$ is
meager, and in particular has the Baire property.
\end{proof}

The same statement holds on any countably infinite set after choosing an
enumeration. To refute the Baire property, it therefore suffices to show
that every finite-block partition has an infinite union of blocks
belonging to the ideal. Our next lemma proves a somewhat stronger
assertion.

\section{Block-support ideals and the fusion argument}
\label{sec:fusion}

\subsection{Block support and countably many transformations}
Let $X$ be a countably infinite set and let $(B_n)_{n\in\omega}$ be a
partition of $X$ into nonempty finite blocks. Define
\begin{equation}
\label{eq:support}
              \supp(A)=\{n:A\cap B_n\ne\varnothing\}.
\end{equation}
The simple block-support ideal is
\[
              \J_\U=\{A\subseteq X:\supp(A)\notin\U\}.
\]
Because $\supp(A\cup D)=\supp(A)\cup\supp(D)$, this is an ideal.
Notice that being small here does not mean having few points in each
block. It means avoiding all points of a $\U$-large collection of
blocks.

For a slightly more general construction, let
$f_i:D_i\to X$, $i\in\omega$, be partial maps with finite fibers, and
assume the identity map on $X$ occurs in the list. Set
\[
              T_i(A)=f_i[A\cap D_i]
\]
and define
\begin{equation}
\label{eq:generalideal}
 \I=\{A\subseteq X:\supp(T_i(A))\notin\U
                               \text{ for every }i\in\omega\}.
\end{equation}
This is again an ideal: each $T_i$ preserves inclusions and finite
unions, sends finite sets to finite sets, and the identity term ensures
$X\notin\I$. The finite-fiber assumption will be needed in the fusion
argument, not in this elementary ideal verification.

\begin{lemma}[Separated-block fusion]
\label{lem:fusion}
For the ideal~\eqref{eq:generalideal}, let $(Q_j)_{j\in\omega}$ be any
sequence of pairwise disjoint, nonempty finite subsets of $X$. There
exists an infinite $H\subseteq\omega$ for which
\[
                        \bigcup_{j\in H} Q_j\in\I.
\]
The $Q_j$ need not cover $X$.
\end{lemma}
\begin{proof}
We first choose increasing indices $j_0<j_1<\cdots$ with the following
finite-stage invariant:
\begin{equation}
\label{eq:fusioninvariant}
 \supp(T_i(Q_{j_n}))\cap\supp(T_\ell(Q_{j_r}))=\varnothing
 \quad\text{if }r<n\text{ and }i,\ell\le n.
\end{equation}
Suppose $j_0,\ldots,j_{n-1}$ have been selected. The set of previously
used block indices, with transformations through stage $n$, is
\[
 H_n=\bigcup_{r<n}\ \bigcup_{\ell\le n}
                         \supp(T_\ell(Q_{j_r})).
\]
It is finite. Consequently $C_n=\bigcup_{b\in H_n}B_b$ is finite, and so
is
\[
                  F_n=\bigcup_{i\le n} f_i^{-1}(C_n).
\]
Preimages here are taken in the domains $D_i$. Finiteness follows from
the finite-fiber hypothesis. Only finitely many of the disjoint sets
$Q_j$ meet $F_n$. Choose $j_n>j_{n-1}$ outside those exceptions. Then
$T_i(Q_{j_n})\cap C_n=\varnothing$ for every $i\le n$, which is exactly
what is needed for~\eqref{eq:fusioninvariant}.

Separate the selected sets into two unions:
\[
 C_0=\bigcup_{n\text{ even}}Q_{j_n},
 \qquad
 C_1=\bigcup_{n\text{ odd}}Q_{j_n}.
\]
For fixed $i,\ell$, the intersection
\begin{equation}
\label{eq:crossfinite}
                  \supp(T_i(C_0))\cap\supp(T_\ell(C_1))
\end{equation}
is finite. Indeed, if a block index lies in this intersection, it is
witnessed by an even selected stage $n$ and an odd selected stage $r$.
These stages are distinct. If $\max(n,r)\ge\max(i,\ell)$,
condition~\eqref{eq:fusioninvariant}, with the larger stage taken last,
forbids the intersection. All remaining witnesses come from finitely
many early selected sets, and their transformed block supports are
finite.

If neither $C_0$ nor $C_1$ belonged to $\I$, there would exist indices
$i$ and $\ell$ with
\[
                 \supp(T_i(C_0))\in\U,
       \qquad    \supp(T_\ell(C_1))\in\U.
\]
Their intersection would belong to $\U$, contradicting its finiteness
in~\eqref{eq:crossfinite} and the freeness of $\U$. Hence at least one
of $C_0,C_1$ belongs to $\I$. Its defining parity supplies the required
infinite index set $H$.
\end{proof}

\begin{corollary}
\label{cor:nonbp}
The ideal~\eqref{eq:generalideal} does not have the Baire property and
is tall. For every $D\in\I^+$, its restriction
\[
                 \I\restr D=\{A\subseteq D:A\in\I\}
\]
does not have the Baire property as a subset of $2^D$.
\end{corollary}
\begin{proof}
Apply Lemma~\ref{lem:fusion} to each finite-block partition supplied by
a putative application of Lemma~\ref{lem:criterion}. The resulting
infinite union contradicts that criterion, so $\I$ cannot have the
Baire property.

For tallness, enumerate an arbitrary infinite $D$ by distinct points
$d_j$ and apply the fusion lemma to $Q_j=\{d_j\}$. It produces an
infinite subset of $D$ belonging to $\I$.

Finally, if $D$ is positive, $\I\restr D$ is a proper ideal on $D$.
Any finite-block partition of $D$ is an eligible sequence $(Q_j)$ in the
fusion lemma, even though it does not cover $X$. The same contradiction
to the category criterion in $2^D$ proves the last assertion.
\end{proof}

\subsection{Taking the translation core}
Return to $X=\omega$. Use the partial maps $f_k(n)=n+k$ on the domains
where $n+k\ge0$, enumerated over $k\in\mathbb Z$. They are injective,
and their induced transformations are the truncated shifts.
Define
\begin{equation}
\label{eq:core}
 \boxed{\displaystyle
 \I_\U=\{A\subseteq\omega:
          \supp(\sh{k}{A})\notin\U\text{ for every }k\in\mathbb Z\}.}
\end{equation}
This is the translation core of $\J_\U$.

\begin{proposition}
\label{prop:core}
For every partition of $\omega$ into nonempty finite blocks,
$\I_\U$ is a proper, tall, translation-invariant ideal without the
Baire property. Its restriction to each positive set also lacks the
Baire property.
\end{proposition}
\begin{proof}
All assertions except translation invariance follow from
Corollary~\ref{cor:nonbp}. For invariance, suppose $A\in\I_\U$ and fix
$h\in\mathbb Z$. For each $k$,
\[
       \supp(\sh{k}{\sh{h}{A}})
                   \eqfin\supp(\sh{k+h}{A}).
\]
The right side does not belong to $\U$, by~\eqref{eq:core}, and finite
changes do not affect ultrafilter membership. Thus
$\sh{h}{A}\in\I_\U$. Conversely, if $\sh{h}{A}\in\I_\U$, apply the
already established closure to $-h$. Equation~\eqref{eq:composition}
and invariance of ideals under finite changes recover $A\in\I_\U$.
\end{proof}

\begin{remark}[Why the category proof needs care]
An embedding of a free ultrafilter into a trace of an ideal does not, by
itself, prove that the whole ideal lacks the Baire property. The Baire
property is not preserved by preimages under arbitrary continuous maps.
For instance, every subset of a closed nowhere dense copy of Cantor
space is meager in the ambient Cantor space, even when its pullback
under the embedding has no Baire property. The fusion argument avoids
this pitfall. It also uses a single parity choice for \emph{all}
transformations, not a potentially different parity for each shift.
\end{remark}

The preservation of non-Baire behavior under translation cores is
already part of the surrounding theory: see~\cite[Proposition~3.8]{FT}
and the countable-intersection theorem attributed there to
Plewik~\cite{Plewik}. The point of the direct argument is to verify the
specific ideal used here without leaving the preservation step implicit.

\section{An explicit separated binary tree}
\label{sec:tree}

\subsection{Blocks and labels}
We now choose the block partition that supplies continuum many positive
sets. For $n\ge0$ put
\begin{equation}
\label{eq:blocks}
\begin{aligned}
 d_n&=2^n,\qquad L_n=4^n+2^{n+1},\\
 b_n&=\sum_{j<n}L_j=\frac{4^n-1}{3}+2(2^n-1),\\
 B_n&=[b_n,b_{n+1})\cap\omega.
\end{aligned}
\end{equation}
Thus $b_0=0$, the intervals partition $\omega$, and each has length
$L_n$. Let $2^{<\omega}$ denote the finite binary words. If $s$ has
length $n$, let $\val(s)$ be its usual binary value between $0$ and
$2^n-1$, with $\val(\varnothing)=0$. Assign the integer label
\begin{equation}
\label{eq:label}
                  p_s=b_n+2^n\bigl(1+\val(s)\bigr).
\end{equation}
The labels at level $n$ form the finite set
\[
                   T_n=\{p_s:|s|=n\}.
\]
They lie inside $B_n$, and
\begin{equation}
\label{eq:margins}
 \min T_n=b_n+2^n,\qquad
 \max T_n=b_n+4^n,\qquad
 b_{n+1}-\max T_n=2^{n+1}.
\end{equation}
Consecutive labels at level $n$ are $2^n$ apart. The gap between the last
label of level $n$ and the first label of level $n+1$ is
\begin{equation}
\label{eq:levelgap}
                     \min T_{n+1}-\max T_n=2^{n+2}.
\end{equation}

\begin{table}[htbp]
\centering
\small
\begin{tabular}{@{}rrrrl@{}}
\toprule
$n$ & $b_n$ & $b_{n+1}$ & $|T_n|$ & Labels in $T_n$\\
\midrule
0 & 0 & 3 & 1 & $1$\\
1 & 3 & 11 & 2 & $5,7$\\
2 & 11 & 35 & 4 & $15,19,23,27$\\
3 & 35 & 115 & 8 & $43,51,59,67,75,83,91,99$\\
4 & 115 & 403 & 16 & $131,147,\ldots,371$\\
\bottomrule
\end{tabular}
\caption{The first levels. Block endpoints are half-open; the labels
stay away from both endpoints by margins tending to infinity.}
\label{tab:levels}
\end{table}

\begin{lemma}[Global separation]
\label{lem:separation}
If $s\ne t$ are finite binary words and $m=\max(|s|,|t|)$, then
\begin{equation}
\label{eq:separation}
                            |p_s-p_t|\ge2^m.
\end{equation}
\end{lemma}
\begin{proof}
If both words have length $m$, their distinct binary values and
formula~\eqref{eq:label} give the bound. If the levels differ, the label
at the higher level is to the right of every label at a lower level.
For $m\ge1$, even its distance from the last label of the immediately
preceding level is
\[
                         \min T_m-\max T_{m-1}=2^{m+1}.
\]
All earlier labels are still farther away. This is stronger than the
claimed bound.
\end{proof}

\subsection{Branches, positivity, and translated intersections}
For $x\in2^\omega$, write $x\restr n$ for its prefix of length $n$ and
set
\begin{equation}
\label{eq:branch}
                       A_x=\{p_{x\restr n}:n\ge0\}.
\end{equation}
Every branch has exactly one point in every $B_n$. Therefore
\begin{equation}
\label{eq:positive}
                         \supp(A_x)=\omega\in\U,
                 \qquad A_x\notin\I_\U.
\end{equation}
The identity shift alone witnesses positivity.

\begin{proposition}[Quantitative translated intersections]
\label{prop:quantitative}
Let $x,y\in2^\omega$.
If $k\ne0$, define
\[
                   N(k)=\lfloor\log_2|k|\rfloor+1.
\]
Then, without any requirement that $x\ne y$,
\begin{equation}
\label{eq:cutoff}
 A_x\cap\sh{k}{A_y}\subseteq[0,b_{N(k)}),
 \qquad
 \bigl|A_x\cap\sh{k}{A_y}\bigr|\le N(k).
\end{equation}
If $x\ne y$ and $d=\min\{j:x(j)\ne y(j)\}$, then
\begin{equation}
\label{eq:commonprefix}
                  |A_x\cap A_y|=d+1.
\end{equation}
\end{proposition}
\begin{proof}
For $k\ne0$, an intersection point gives
$p_{x\restr n}=p_{y\restr r}+k$. The two node labels are distinct.
By Lemma~\ref{lem:separation},
\[
                           |k|\ge2^{\max(n,r)}.
\]
Since $2^{N(k)}>|k|$, both $n$ and $r$ are smaller than $N(k)$. Thus the
intersection is contained in the first $N(k)$ levels, all of which lie
below $b_{N(k)}$. There is at most one point of $A_x$ in each such level,
proving~\eqref{eq:cutoff}.

For $k=0$, equality of node labels forces equality of the words, since
distinct levels are in disjoint blocks and labels within a level are
distinct. Two different infinite branches have the same prefixes
exactly at lengths $0,1,\ldots,d$. This includes the empty prefix and
explains the $+1$ in~\eqref{eq:commonprefix}.
\end{proof}

The cutoff is a uniform bound for each fixed nonzero shift, independent
of both branches. It is not asserted to be optimal. The dependence on
$k$ is essential, and the zero shift is treated separately because two
branches can share an arbitrarily long finite prefix.

\begin{proposition}[The perfect family]
\label{prop:perfect}
The map $\Phi(x)=A_x$ is a continuous injection from $2^\omega$ into
$2^\omega$. Its image is compact and perfect and has cardinality
$\cfr$.
\end{proposition}
\begin{proof}
For a coordinate $q$ that is not a node label, membership
$q\in A_x$ is constantly false. If $q=p_s$, it is equivalent to $s$
being a prefix of $x$, a clopen condition on $x$. Hence every coordinate
map is continuous, so $\Phi$ is continuous in the product topology.
Distinct branches differ at some prefix, and their sets consequently
differ at the corresponding label; thus $\Phi$ is injective.

A continuous injection from the compact space $2^\omega$ into the
Hausdorff space $2^\omega$ is a homeomorphism onto its image. The image
is therefore compact, has no isolated points, and has the cardinality
of Cantor space.
\end{proof}

\begin{proof}[Proof of Theorem~\ref{thm:main}]
Use the partition~\eqref{eq:blocks} in~\eqref{eq:core}.
Proposition~\ref{prop:core} establishes all assertions about the ideal,
including tallness and the category of its positive restrictions.
Equation~\eqref{eq:positive} establishes positivity of every branch.
Proposition~\ref{prop:quantitative} gives the finite-intersection
properties, and Proposition~\ref{prop:perfect} gives the required
perfect parametrization. Since $\Fin\subseteq\I_\U$, the family is
in particular $\I_\U$-TAD. This proves the stated affirmative answer to
Problem~\ref{prob:FT}.
\end{proof}

\section{The ultrafilter trace and the quotient Boolean algebra}
\label{sec:trace}

The increasing margins in~\eqref{eq:margins} have another purpose:
a fixed translation eventually keeps every selected point in its own
block. This identifies the ideal on every branch exactly.

For $S\subseteq\omega$, define
\[
                   A_x[S]=\{p_{x\restr n}:n\in S\}.
\]
\begin{proposition}[Exact trace formula]
\label{prop:trace}
For every $x\in2^\omega$, $S\subseteq\omega$, and $k\in\mathbb Z$,
\begin{equation}
\label{eq:tracesupport}
                         \supp(\sh{k}{A_x[S]})\eqfin S.
\end{equation}
Consequently,
\begin{equation}
\label{eq:trace}
                         A_x[S]\in\I_\U
                              \quad\Longleftrightarrow\quad S\notin\U.
\end{equation}
\end{proposition}
\begin{proof}
Whenever $2^n>|k|$, the margin estimates show that
$p_{x\restr n}+k$ is still in $B_n$. Thus all sufficiently late selected
points retain their block index. The finitely many early selected
points contribute only finitely many block indices before and after
translation, proving~\eqref{eq:tracesupport}.

If $S\notin\U$, every support in~\eqref{eq:tracesupport} is outside
$\U$, so $A_x[S]\in\I_\U$. If $S\in\U$, the identity shift has support
exactly $S$, and hence $A_x[S]\notin\I_\U$.
\end{proof}

Thus $n\mapsto p_{x\restr n}$ transports the dual ideal of $\U$ to
$\I_\U\restr A_x$. In particular, this trace is a maximal ideal on
$A_x$. Notice that this does not make the whole ideal maximal: there
are many different positive branches.

Recall that $\Pow(\omega)/\I_\U$ identifies $A$ and $D$ when
$A\mathbin\triangle D\in\I_\U$. An \emph{atom} in a Boolean algebra is
a nonzero element with no strictly intermediate element between it and
zero.

\begin{corollary}
\label{cor:atoms}
Each $[A_x]$ is an atom in $\Pow(\omega)/\I_\U$. The elements $[A_x]$
are distinct and pairwise disjoint, so the quotient has at least
$\cfr$ atoms. Distinct free ultrafilters give distinct ideals in the
construction.
\end{corollary}
\begin{proof}
Every subset of $A_x$ has the form $A_x[S]$. Exactly one of $S$ and its
complement belongs to $\U$. By~\eqref{eq:trace}, either $A_x[S]$ is zero
in the quotient, or its complement relative to $A_x$ is zero there.
This proves atomicity. More generally, an element below $[A_x]$ can be
represented by a subset of $A_x$, by intersecting a representative with
$A_x$.

For $x\ne y$, the intersection $A_x\cap A_y$ is finite, so the two
nonzero atoms are disjoint and distinct. Finally, if $\U\ne\mathcal V$,
choose $S$ belonging to one ultrafilter and not the other. Formula
\eqref{eq:trace}, using any one fixed branch, distinguishes
$\I_\U$ from $\I_{\mathcal V}$.
\end{proof}

The occurrence of continuum many atoms is compatible with the
countability of the ground set. These atoms live in a quotient by an
ideal; the representing sets are not literally disjoint, only disjoint
modulo the ideal. No claim is made that these are all the atoms.

\section{Abstract densities with this exact null ideal}
\label{sec:density}

\subsection{The meaning of abstract upper density}
An \emph{abstract upper density} is a map
$\delta:\Pow(\omega)\to[0,1]$ such that $\delta(\omega)=1$, finite sets
have value zero, $\delta$ is monotone, and
\[
                       \delta(A\cup D)\le\delta(A)+\delta(D).
\]
It is \emph{rich} when its range is the whole interval $[0,1]$, and
translation invariant when $\delta(\sh{k}{A})=\delta(A)$ for every $k$.
This is the framework of~\cite{DNJ,FT}. It does not require agreement
with ordinary asymptotic density, finite additivity, or a dilation
identity. The distinction matters for our example.

\subsection{A maxitive realization and prescribed branch values}
The next argument is an application of the established maximal-TAD
construction~\cite[Theorem~3.1]{FT}, related to the maximal-family method
of~\cite{DNJ}. We give the formula and verify it, rather than claiming
the general mechanism as a new theorem.

\begin{theorem}
\label{thm:density}
Let $w:2^\omega\to(0,1]$ be any surjection. There is a translation-invariant
rich abstract upper density $\delta$ with
\begin{equation}
\label{eq:exactnull}
          \{A:\delta(A)=0\}=\I_\U,
        \qquad \delta(A_x)=w(x)\quad(x\in2^\omega).
\end{equation}
It can be chosen \emph{maxitive}:
\begin{equation}
\label{eq:maxitive}
                    \delta(A\cup D)=\max\{\delta(A),\delta(D)\}.
\end{equation}
\end{theorem}
\begin{proof}
Extend $\{A_x:x\in2^\omega\}$ to a maximal $\I_\U$-TAD family
$\M\subseteq\I_\U^+$. Zorn's lemma applies because the union of a chain
of extending TAD families is again such a family: any two of its
members already lie together in one element of the chain.

Define a weight $a:\M\to(0,1]$ by $a(A_x)=w(x)$ and $a(C)=1$ on any
remaining members. All weights are strictly positive. For
$A\subseteq\omega$, let
\[
 W(A)=\{C\in\M:\text{there is }k\in\mathbb Z
                         \text{ with }C\cap\sh{k}{A}\notin\I_\U\},
\]
and set
\begin{equation}
\label{eq:densityformula}
                    \delta(A)=\sup\{a(C):C\in W(A)\},
                        \qquad\sup\varnothing=0.
\end{equation}
We verify each required property.

If $A\subseteq D$, then $W(A)\subseteq W(D)$, so $\delta$ is monotone.
A finite $A$ has no witnesses, so its value is zero. Every member of
$\M$ witnesses $\omega$, and $w$ takes the value~1; hence
$\delta(\omega)=1$.

For any fixed $C$ and $k$, the intersection with a shifted union is
\[
 C\cap\sh{k}{A\cup D}
       =\bigl(C\cap\sh{k}{A}\bigr)
                         \cup\bigl(C\cap\sh{k}{D}\bigr).
\]
A union belongs to an ideal exactly when both constituent sets do.
It follows that $W(A\cup D)=W(A)\cup W(D)$.
Taking suprema proves~\eqref{eq:maxitive}, which implies subadditivity.

For translation invariance, $\sh{k}{\sh{h}{A}}$ differs finitely from
$\sh{k+h}{A}$, and adding or removing finitely many points does not
change whether an intersection is positive. As $k$ ranges over
$\mathbb Z$, so does $k+h$. Thus $W(\sh{h}{A})=W(A)$.

If $A\in\I_\U$, every shift and every intersection of a shift with a
member of $\M$ belongs to $\I_\U$, so $W(A)=\varnothing$. Conversely,
suppose $A\notin\I_\U$ and $W(A)=\varnothing$. Then $A$ is translated
almost disjoint from every member of $\M$. The condition is symmetric:
shifting $C\cap\sh{k}{A}$ back by $-k$ gives
$A\cap\sh{-k}{C}$ modulo a finite set. Also $A\notin\M$, since otherwise
it witnesses itself with $k=0$. Therefore adjoining $A$ would strictly
enlarge $\M$, a contradiction. Hence every positive $A$ has a witness,
and its density is strictly positive. This proves the exact-null-ideal
assertion.

Finally, $W(A_x)=\{A_x\}$. The inclusion of $A_x$ follows from
positivity; no other member of the TAD family can witness it. Therefore
$\delta(A_x)=a(A_x)=w(x)$. Surjectivity of $w$ supplies every value in
$(0,1]$, and the empty set supplies zero.
\end{proof}

A surjection $w$ can be chosen very concretely: take
$r(x)=\sum_{j\ge0}x(j)2^{-j-1}$, and replace its value~0 by~1. This has
range $(0,1]$. Neither continuity nor measurability of the weights is
required.

\subsection{Why finite additivity is impossible}
\begin{proposition}
\label{prop:noadditive}
There is no finitely additive probability measure on all subsets of
$\omega$ whose null ideal is exactly $\I_\U$.
\end{proposition}
\begin{proof}
Suppose $\mu$ were such a measure. Every $A_x$ would have positive
measure, and all intersections $A_x\cap A_y$ with $x\ne y$ would have
measure zero. For any finite list of distinct branches, successively
remove the intersections with earlier branches. This makes the list
literally disjoint without changing any of its measures. Therefore
\[
                       \sum_{j=1}^r\mu(A_{x_j})\le1.
\]
For each positive integer $m$, at most $m$ branches can have measure
at least $1/m$. But every positive real number is at least $1/m$ for
some $m$. Thus only countably many branches could have positive
measure, contradicting their continuum cardinality.
\end{proof}

This is the usual countable-chain-condition obstruction for a measure
algebra, applied directly. In particular, the nonadditivity of the
maxitive density is not an accidental weakness of the formula.
Moreover, any function with exact zero set $\I_\U$ cannot be Baire
measurable: a Baire-measurable real-valued function has a zero set with
the Baire property, since that zero set is a countable intersection of
preimages of open intervals.

\section{Comparison with ordinary density}
\label{sec:banach}

For clarity, define upper Banach density here by
\[
 d^*(E)=\limsup_{L\to\infty}\ \sup_{a\in\omega}
                   \frac{|E\cap[a,a+L)|}{L},
              \qquad L\in\mathbb N_{>0}.
\]
We need only its defining estimates, not any structural theorem about
this density. Let $\mathcal Z_B=\{E:d^*(E)=0\}$.

\begin{proposition}
\label{prop:incomparable}
The ideals $\I_\U$ and $\mathcal Z_B$ are incomparable under inclusion.
Every branch $A_x$ has upper Banach density zero but is
$\I_\U$-positive. Conversely, $\I_\U$ contains a set of upper Banach
density one.
\end{proposition}
\begin{proof}
Let $T=\bigcup_n T_n$ be the set of all tree labels. Its successive gaps
tend to infinity by~\eqref{eq:margins} and~\eqref{eq:levelgap}. Given a
positive integer $M$, remove a finite initial set $F\subseteq T$ so
that consecutive points of $T\setminus F$ are at least $M$ apart.
Any interval of length $L$ contains at most $L/M+1$ such points, so
\[
       \sup_a\frac{|T\cap[a,a+L)|}{L}
                       \le\frac1M+\frac{|F|+1}{L}.
\]
First let $L\to\infty$, then $M\to\infty$. This gives $d^*(T)=0$ and
hence $d^*(A_x)=0$. Positivity was proved in~\eqref{eq:positive}.

For the other direction, the four residue classes modulo~4 partition
$\omega$, so exactly one of them, say $R_r$, belongs to $\U$. Let
$S=R_{r+2\bmod4}$ and put
\[
                         E=\bigcup_{n\in S}B_n.
\]
Fix $k\in\mathbb Z$. Since $L_n\to\infty$, all sufficiently late
blocks and their neighbors have length larger than $|k|$.
Translating a late block by $k$ can therefore meet only itself and its
immediate neighbors. Up to finitely many block indices,
\[
                 \supp(\sh{k}{E})\subseteq S+\{-1,0,1\}.
\]
The set on the right misses $R_r$. Consequently the support has finite
intersection with the member $R_r$ of $\U$ and cannot belong to $\U$.
This is true for every $k$, so $E\in\I_\U$.

On the other hand, $E$ contains the full intervals $B_n$ for infinitely
many $n$, with lengths tending to infinity. For each fixed $L$, some
such interval has length at least $L$; it contains an interval of length
$L$ entirely inside $E$. Thus the supremum in the definition of $d^*$ is
1 for every $L$, and $d^*(E)=1$.
\end{proof}

It follows in particular that the density in
Theorem~\ref{thm:density} does not extend ordinary asymptotic density:
it assigns a positive value to sets $A_x$ whose asymptotic density is
zero. Nor can it dominate upper Banach density, because it vanishes on
the set $E$ just constructed. These observations locate the result
inside the \emph{abstract} density framework and prevent an unintended
stronger interpretation.

\section{Extension to countable permutation actions}
\label{sec:group}

The mechanism is not specific to integer addition. Countability of the
transformations and finiteness of their preimages of finite sets are
the ingredients used in the category proof. The tree can also be built
by finite avoidance rather than by the explicit numerical formula.

\begin{theorem}
\label{thm:group}
Let a countable group $G$ act by permutations on a countably infinite
set $X$. In ZFC there is a proper, tall, $G$-invariant ideal $\I$ on $X$
without the Baire property and a perfect family
$\{A_x:x\in2^\omega\}\subseteq\I^+$ such that
\[
                         A_x\cap gA_y\text{ is finite}
               \quad(x\ne y,\ g\in G).
\]
The restriction of $\I$ to each positive set also lacks the Baire
property. If the action is free, $A_x\cap gA_x$ is finite for every
$g\ne e$ as well.
\end{theorem}
\begin{proof}
Choose an enumeration of $X$ and increasing finite symmetric sets
$G_n\subseteq G$ with $e\in G_n$ and $\bigcup_nG_n=G$.
We recursively build finite blocks $B_n$ and distinct labels $p_s$ for
all binary words $s$, putting the labels of words of length $n$ into
$B_n$.

At stage $n$, the previously constructed blocks form a finite set.
Let $e_n$ be the first point of the enumeration not yet used in a block.
Reserve it as a filler. Choose the $2^n$ labels of the new level one at
a time. A new label must avoid all previous blocks, the filler $e_n$,
the already chosen labels at this level, and
\[
              \{g p_t:g\in G_n,\ t\text{ labels a previously
                                      chosen tree node}\}.
\]
Only finitely many points are forbidden at each choice, so the infinite
set $X$ always supplies another point. Symmetry of $G_n$ ensures that
avoiding $p_s=g p_t$ also avoids $p_t=g p_s$. When the level has been
chosen, let $B_n$ be its labels together with $e_n$.

The blocks are pairwise disjoint and finite. They cover $X$: at every
stage the first unused point is placed in a block, so every fixed point
of the enumeration is eventually used. For distinct tree nodes $s,t$,
the construction guarantees
\begin{equation}
\label{eq:groupseparation}
 p_s\ne g p_t
 \quad\text{whenever }g\in G_N
                    \text{ and }\max(|s|,|t|)\ge N.
\end{equation}
Indeed, at the later of the two selection times the element $g$ and its
inverse belong to the current finite symmetric set, and the relevant
point is forbidden. This also covers two distinct nodes chosen at the
same level.

Fix a free ultrafilter $\U$ on $\omega$, use block support for this
partition, and put
\[
            \I=\{A\subseteq X:\supp(gA)\notin\U
                                           \text{ for every }g\in G\}.
\]
It is proper because the identity transformation sends $X$ to $X$.
The fusion lemma and Corollary~\ref{cor:nonbp} apply to the permutation
maps, giving the category and tallness assertions. For $h\in G$,
reindexing $g$ by $gh$ proves $hA\in\I$ if and only if $A\in\I$.

Define $A_x=\{p_{x\restr n}:n\ge0\}$. It has exactly one point in each
block and is positive. For $g\in G_N$, an equality between a point of
$A_x$ and a point of $gA_y$ either involves the same tree node or,
by~\eqref{eq:groupseparation}, involves two nodes at levels below $N$.
Distinct branches have only finitely many common nodes, and there are
only finitely many nodes at levels below $N$. This proves the required
finite intersection. In a free action a nonidentity $g$ fixes no point,
so the same-node possibility disappears even for $x=y$.

Finally, membership of a label $p_s$ in $A_x$ is exactly the clopen
prefix condition, while a filler point belongs to no branch. The same
compactness argument as in Proposition~\ref{prop:perfect} gives a
perfect family.
\end{proof}

For a nonfree action the final self-intersection assertion need not
hold: the identity action fixes every branch. This is why the theorem
separates distinct-branch intersections from self-intersections. The
main result on $\omega$ used explicit truncated shifts; the theorem
above concerns genuine permutation actions and does not silently
identify the two settings.

\section{Verification, reproducibility, and proof audit}
\label{sec:verification}

\subsection{What the program checks}
The accompanying \texttt{verify.py} uses only the Python standard
library and exact integer arithmetic. It contains explicit checks that
remain active even under Python's optimization flag. A failed check
raises an error and terminates with a nonzero exit status.

The recorded run completed \textbf{13,012,145 successful checks} with
Python~3.13.5. These include the replay of the random draw, the block
boundary formula, uniqueness and placement of all 8,191 tree labels
through level~12, their endpoint margins, and adjacent-label separation.
For every nonzero integer shift from $-512$ to $512$, the program scans
all those labels for collisions and checks the theoretical cutoff.

All 256 binary branch prefixes of length~8 are also tested. Their
zero-shift intersections are checked against the exact common-prefix
formula. Every ordered pair, including identical prefixes, is checked
for every nonzero shift from $-31$ to $31$. Both the cardinality bound
and the location bound are verified. A further set of tests checks
same-block transport under shifts from $-256$ to $256$ whenever the
margin hypothesis applies.

Finally, eight finite fusion constructions are run through 24 stages
each, using separate reproducible seeds and auxiliary blocks of lengths
$3,4,5,\ldots$. The program checks every finite-stage separation
condition, then independently verifies that each even--odd transformed
support overlap is confined to the predicted early stages. The
auxiliary blocks keep the test inexpensive; the written fusion lemma
applies to \emph{every} finite-block partition, including the main
exponentially growing one.

\subsection{Files and reproducibility}
\begin{center}
\small
\begin{tabular}{@{}p{0.39\textwidth}p{0.55\textwidth}@{}}
\toprule
File & Purpose\\
\midrule
\texttt{article.tex}, \texttt{article.pdf} & Source and compiled article.\\
\texttt{verify.py} & Exact finite verification program.\\
\texttt{random\_selection.json} & Original area pool, seed, and selected index.\\
\texttt{verification\_results.json} & Parameters, check counts, and limitations.\\
\texttt{verification\_report.txt} & Readable summary of the recorded run.\\
\texttt{tree\_levels.csv} & Exact block and label data through level~12.\\
\texttt{shift\_collisions.csv} & Collision counts for the finite tree and tested shifts.\\
\texttt{fusion\_certificates.json} & Selected finite-block subsequences in the fusion tests.\\
\texttt{literature\_search.md} & Sources consulted and search limitations.\\
\texttt{README.md}, \texttt{Makefile} & Build and execution instructions.\\
\bottomrule
\end{tabular}
\end{center}

To rerun the checks, execute \texttt{python3 verify.py}. To rebuild the
PDF, execute \texttt{latexmk -pdf article.tex}, or run
\texttt{pdflatex article.tex} twice. No network access is needed for
these steps once the archive has been obtained. The compiled article
uses ordinary LaTeX packages; its bibliography is embedded in the
source.

\subsection{The infinitary checks are mathematical, not computational}
No finite object is substituted for a free ultrafilter. On a finite
set every ultrafilter is principal, which would invalidate the
finite-intersection contradiction in the fusion lemma. Accordingly,
the program does not claim to query $\U$, test ideal membership, or
verify failure of the Baire property.

The main logical dependencies can be audited as follows. First, the
category criterion produces a finite-block obstruction from the
assumption of the Baire property. Second, the fusion lemma defeats
every such obstruction using finite fibers and one fixed free
ultrafilter. Third, the orbit-core definition enforces invariance;
truncated-shift compositions are used only modulo finite sets. Fourth,
the branch sets are positive because they hit every block, not because
of an assumption about their density. Fifth, global separation treats
collisions between different levels as well as within a level. Finally,
continuity and injectivity, rather than a bare cardinality assertion,
produce the perfect family.

The proof of the density corollary has two further important points:
all weights are strictly positive, so a nonempty witness set cannot
have supremum zero; and maximality is invoked for exactness of the
null ideal, not merely for the availability of many prescribed values.
The group extension uses symmetric finite sets of group elements so
that avoidance works in both directions.

\section{Conclusion and remaining scope}
\label{sec:conclusion}

The construction separates two tasks that initially look in tension.
Ultrafilter block supports make every full block selector positive,
while a translation core makes the ideal invariant. Neither operation
requires selectors to be close together. By leaving enough room inside
each block for widely separated binary-tree labels, one obtains
continuum many positive selectors with only finite translated
intersections. The fusion lemma then supplies the missing category
argument for the core.

Thus the mathematical argument presented here gives an affirmative
construction for the stated Question~1, and in fact a perfect-family
version with explicit bounds. The further consequences show that the
example is tall, has maximal ultrafilter traces on its branches, admits
a rich maxitive abstract density with its exact null ideal, and has an
analogue for countable permutation actions.

There are definite limits. The construction does not settle whether
every translation-invariant ideal is the null ideal of a rich invariant
abstract density. It does not construct such a density for a
translation-maximal ideal. In fact $\I_\U$ is not translation maximal:
fix distinct branches $A_x,A_y$ and form the invariant ideal generated
by $\I_\U$ together with all shifts of $A_x$. If that ideal contained
$\omega$, then $A_y$ would be covered by an ideal-small set and finitely
many shifts of $A_x$. Its intersections with those shifts are finite,
forcing $A_y\in\I_\U$, a contradiction. The generated ideal is
therefore proper and strictly larger.

Nor is a definable or computable replacement for the free ultrafilter
provided. A proof-assistant formalization and an independent
mathematical review remain distinct verification tasks, and the
literature-priority question remains separate from the correctness of
the explicit argument.

\clearpage
\appendix
\section{Random-area selection record}
\label{app:random}

The original decimal seed was
\begin{center}
\texttt{209442498745758881649029240816391254350}.
\end{center}
It was generated using \texttt{secrets.randbits(128)}. A local instance
of \texttt{random.Random(seed)} then selected
\texttt{randrange(24)}, obtaining \textbf{22}. The recorded UTC time was
September 20, 2026, at 21:02:49.492499 UTC. The complete pool, in its
original order, was:

\begin{center}
\small
\begin{longtable}{@{}rl@{}}
\toprule
Index & Area\\
\midrule
\endhead
0 & Analytic number theory\\
1 & Algebraic number theory\\
2 & Finite group theory\\
3 & Commutative algebra\\
4 & Discrete geometry\\
5 & Extremal graph theory\\
6 & Probability and random processes\\
7 & Dynamical systems\\
8 & Approximation theory and special functions\\
9 & Matrix analysis and inequalities\\
10 & Topology and knot theory\\
11 & Formal languages and automata\\
12 & Order theory and lattice theory\\
13 & Enumerative combinatorics\\
14 & Additive combinatorics\\
15 & Optimization and convex analysis\\
16 & Coding theory\\
17 & Algebraic geometry\\
18 & Functional equations\\
19 & Combinatorial game theory\\
20 & Spectral graph theory\\
21 & Numerical analysis\\
22 & \textbf{Set theory}\\
23 & Mathematical logic and computability\\
\bottomrule
\end{longtable}
\end{center}

The selected mathematical area was random; the subsequent choice of
Problem~\ref{prob:FT} within that area was based on its precise published
formulation and the availability of a plausible constructive approach.
Random selection does not guarantee that another researcher has not
chosen the same problem.

\section{A compact proof map}
\label{app:proofmap}
For a reader checking the central claim independently, the whole
argument can be reduced to the following sequence of assertions.

Start with any free $\U$ and the blocks
\[
 B_n=\left[\frac{4^n-1}{3}+2(2^n-1),\,
              \frac{4^{n+1}-1}{3}+2(2^{n+1}-1)\right)\cap\omega.
\]
Let $\supp(A)=\{n:A\cap B_n\ne\varnothing\}$ and impose
$\supp(\sh{k}{A})\notin\U$ simultaneously for every integer $k$.
The resulting family is an ideal, is proper at the identity shift,
and is invariant by composition modulo finite sets.

To test the category assertion, begin with an arbitrary finite-block
partition $(Q_j)$. Select a subsequence so that the first $n$ translated
supports of the new block avoid the first $n$ translated supports of
all preceding selected blocks. This is possible because a finite union
of the $B_n$ is finite and each translation has finite preimages. The
even and odd selected unions have finite cross-intersection of their
block supports under every pair of fixed shifts. They cannot both
have a shifted support in a free ultrafilter. One therefore belongs to
the ideal and contains infinitely many $Q_j$, contradicting the
category criterion under a Baire-property assumption.

Finally, put
\[
 A_x=\left\{
      \frac{4^n-1}{3}+2(2^n-1)+2^n(1+\val(x\restr n)):n\ge0
                         \right\}.
\]
These sets are positive full selectors. Distinct node labels with
maximum level $m$ differ by at least $2^m$. A fixed nonzero difference
can therefore occur at only finitely many levels; difference zero
between distinct branches is confined to their common prefixes.
Coordinate membership is a clopen prefix condition, so the family is
perfect. These statements are sufficient for the main theorem without
using any of the density or group-action extensions.

\begin{thebibliography}{9}
\addcontentsline{toc}{section}{References}

\bibitem{FT}
Rafa\l{} Filip\'ow and Jacek Tryba,
\emph{Densities for sets of natural numbers vanishing on a given family},
Journal of Number Theory \textbf{211} (2020), 371--382.
\href{https://doi.org/10.1016/j.jnt.2019.10.015}{doi:10.1016/j.jnt.2019.10.015}.
Preprint: \href{https://arxiv.org/abs/2309.00982}{arXiv:2309.00982v1},
posted September 2, 2023. Question~1 is on preprint page~7.

\bibitem{DNJ}
Mauro Di Nasso and Renling Jin,
\emph{Abstract densities and ideals of sets},
Acta Arithmetica \textbf{185} (2018), no.~4, 301--313.
Preprint: \href{https://arxiv.org/abs/1709.03097}{arXiv:1709.03097}.

\bibitem{Talagrand}
Michel Talagrand,
\emph{Compacts de fonctions mesurables et filtres non mesurables},
Studia Mathematica \textbf{67} (1980), no.~1, 13--43.
Theorem~21; the ideal formulation used here is also stated in
\cite[Theorem~1.7]{FT}.
\href{https://eudml.org/doc/218303}{EuDML record 218303}.

\bibitem{Plewik}
Szymon Plewik,
\emph{Ideals of the second category},
Fundamenta Mathematicae \textbf{138} (1991), no.~1, 23--26.
\href{https://doi.org/10.4064/fm-138-1-23-26}{doi:10.4064/fm-138-1-23-26}.
The countable-intersection result is quoted in
\cite[Theorem~1.9]{FT}; our argument does not depend on invoking it.

\bibitem{KL}
Jonathan M. Keith and Paolo Leonetti,
\emph{Completeness and additive property for submeasures},
\href{https://arxiv.org/abs/2501.13615}{arXiv:2501.13615v2},
March 23, 2026. Consulted for surrounding literature, not used in the
proof.

\bibitem{Path}
Rafa\l{} Filip\'ow and Jacek Tryba,
\emph{Path of pathology},
\href{https://arxiv.org/abs/2501.00503}{arXiv:2501.00503v1},
submitted December 31, 2024. Consulted for surrounding literature, not
used in the proof.

\end{thebibliography}
\end{document}
